\section{Discussion and Future Directions}
\label{sec:discussion}

The simulation design presented in this paper is the first stage of a
broader program: an observationally anchored, progressively extensible
ensemble of resolved multiphase ISM simulations that supports systematic
comparison with \PHANGS\ and statistical inference of ISM physical
parameters. In this section we recapitulate the sampling design
(\autoref{sec:disc_summary}), outline how the block-structured Sobol
construction will be extended to physics parameters
(\autoref{sec:disc_phys}), describe the synthetic observation pipeline
that will close the loop with resolved \PHANGS\ maps
(\autoref{sec:disc_synth}), discuss summary statistics beyond one-point
aperture quantities (\autoref{sec:disc_twopoint}), and frame the
connection to the inference framework developed in the companion paper
(\autoref{sec:disc_paperII}).

\subsection{Summary of the Sampling Design}
\label{sec:disc_summary}

We have presented a \PHANGS-informed simulation design for the
\TIGRESS\ framework. Five environmental parameters---the gas surface
density $\Sgas$, stellar surface density $\Sstar$, stellar scale height
$\Hstar$, orbital frequency $\Omrot$, and shear parameter
$\qshear$---define the design vector $\thetaenv$ and enter \TIGRESS\
as initial and boundary conditions. The sampling proceeds in two steps:
a kernel density estimator (KDE) is fit to the joint \PHANGS\ aperture
distribution of $\thetaenv$ and used to define the marginal quantile
functions, and a scrambled Sobol quasi-random sequence is mapped through
these quantile functions to generate the simulation design.

Two features of this construction are essential for the program as a
whole. First, because the Sobol sequence is deterministic and
hierarchical, the suite can be extended progressively: doubling the
sample size adds new low-discrepancy points without displacing or
re-running any earlier simulation. Second, the molecular fraction
$\fmol$ and SFR surface density $\Ssfr$ are deliberately excluded from
the input design and reserved as validation quantities, to be compared
against emergent \TIGRESS\ outputs.
\CGK{This separation of design fields from validation fields is what
allows the suite to be used both as a calibration set (for \PRFM\
yields) and as a prior for simulation-based inference; flag for review
that this framing is consistent with the abstract.}

\subsection{Extension to Physics Parameter Variations}
\label{sec:disc_phys}

The first suite holds the physics parameters at fiducial
solar-neighborhood values: solar metallicity ($\Zg = 1$), solar
dust-to-gas ratio ($\Zd = 1$), a standard initial mass function (IMF),
and the fiducial \TIGRESS-NCR\ feedback prescription. This choice
isolates the environmental dependence of ISM properties from physics
variations, but it also means that the inferred posteriors from this
suite are conditioned on these fiducial physics assumptions. The
natural next step is to relax this conditioning by introducing an
independent block of simulations that varies a vector of physics
parameters $\thetaphys$.

We anticipate four parameters as the principal targets of this
extension. First, the gas-phase metallicity $\Zg \in [0.1, 3]$ in solar
units, motivated by the \TIGRESS-NCR\ metallicity suite of
\citet{2024ApJ...972...67K}, who showed that $\Ssfr$ scales weakly with
metallicity at fixed environment, approximately as $\Ssfr \propto
\Zg^{0.3}$. \CGK{Verify the exponent and the sign; I am recalling this
from the K24 abstract but please confirm against the paper before this
text is finalized.} Second, the dust-to-gas ratio $\Zd$, which is
strongly but not perfectly correlated with $\Zg$ in observed galaxies
and which controls photoelectric heating, $\HH$ formation, and
radiation transfer independently of gas-phase line cooling. Third, the
upper-mass cutoff of the IMF, which sets the ionizing photon budget and
the supernova rate per unit star formation. Fourth, an overall feedback
yield normalization that re-weights the thermal versus turbulent
partitioning of $\Ytot$ relative to its calibrated \TIGRESS-NCR\ value;
this last parameter is included not because we expect the calibration
to be wrong, but to allow the inference framework to detect
inconsistencies between the simulated and observed yield decompositions.

The block-structured Sobol construction makes this extension natural.
Block~1 is the present environmental suite, a Sobol design of dimension
$d_{\rm env} = 5$ over $\thetaenv$ at fiducial physics. Block~2 will be
a Sobol design of dimension $d_{\rm phys}$ over $\thetaphys$, run at a
representative subset of environmental anchor points selected from
Block~1. The joint suite supports inference over the concatenated
parameter vector $(\thetaenv, \thetaphys)$ by combining the marginal
information from each block under an assumed factorization of the
prior. This mirrors the multi-suite philosophy of the \CAMELS\ project
\citep{2021ApJ...915...71V}, in which one-parameter variation suites,
Latin-hypercube suites, and extension suites jointly cover the
cosmology and astrophysics parameter space at controlled cost.
\CGK{Verify that this is a faithful characterization of the CAMELS suite
structure; the high-level analogy is correct but the suite labels and
their exact roles should be checked.}

\subsection{Synthetic Observation Pipeline}
\label{sec:disc_synth}

A direct comparison between \TIGRESS\ and \PHANGS\ at the level of
integrated aperture quantities, as performed in
\autoref{sec:results}, is a necessary but not sufficient test of the
simulation suite. The same aperture-averaged $\Ssfr$, $\fmol$, $\PDE$,
and $\seff$ can be produced by simulations that differ substantially in
the spatial structure of their gas and young stellar populations. To
test the simulations at the level of resolved morphology, and to
construct simulated observables that can be ingested by the inference
framework of the companion paper, we plan to develop a forward-modeling
pipeline that converts \TIGRESS\ snapshots into mock multi-wavelength
maps matched to the \PHANGS\ resolution and noise properties.

The pipeline will produce five primary tracers. (i)~Mid-infrared dust
emission, modeled from the \TIGRESS\ dust column and ambient radiation
field to produce synthetic JWST/MIRI and \emph{Spitzer}/MIPS $24\,\mu$m
maps. (ii)~Ionized-gas H$\alpha$ emission, computed from the \HII\
region catalog of each \TIGRESS\ snapshot with extinction applied
self-consistently from the simulated dust column. (iii)~Far-ultraviolet
continuum, computed from the young stellar population with the same
extinction treatment. (iv)~\HI\ 21\,cm column, derived from the neutral
atomic phase of \TIGRESS-NCR. (v)~CO line emission, computed either
from a calibrated CO-to-H$_2$ conversion as a function of metallicity
and radiation field, or from direct radiative transfer post-processing
of the \HH\ phase. \CGK{Decide which CO route to adopt; both have a
literature trail and the choice will need to be justified explicitly.}
Each tracer will then be convolved with the appropriate beam, projected
into the observational aperture grid, and resampled to the \PHANGS\
sensitivity and noise level.

This synthetic observation pipeline serves three purposes. It allows
direct comparison of \TIGRESS\ outputs to observed \PHANGS\ maps in the
observed plane, removing the modeling assumptions implicit in
converting observed maps to physical surface densities. It produces the
maps required as inputs to field-level inference (\autoref{sec:disc_paperII}).
And it provides a self-consistent training set for emulators that
predict observable maps directly from $\thetaenv$, bypassing the
intermediate step of physical-parameter estimation.

\subsection{Beyond One-Point Statistics}
\label{sec:disc_twopoint}

The validation and inference targets used in this paper are one-point
aperture summaries: the mean $\Ssfr$, $\fmol$, $\PDE$, $\seff$, and
feedback yield $\Ytot$ within a fixed aperture. These quantities are
informative about the pressure balance and the mean efficiency of
feedback, but they discard the spatial coherence of star formation and
gas distribution that is plausibly the most discriminating signature
between feedback models with similar globally averaged behavior.

A natural first extension is to two-point statistics: the angular power
spectrum or two-point correlation function of $\Sgas$, $\Ssfr$, and
$\PDE$ maps within each aperture or across each simulated patch.
Two-point statistics constrain the spatial coherence length of star
formation, the slope of the turbulent cascade in the cold gas, and the
relative phasing of gas and SFR tracers. \CGK{Mention specific
turbulent-cascade citations here if desired; I have left this generic
because I do not have a high-confidence ADS key in hand for the most
appropriate references.}

A more powerful and increasingly mature class of statistics for
non-Gaussian fields is the wavelet scattering transform
\citep{CITATION_NEEDED_TBD}. Wavelet scattering coefficients are
translation-equivariant, rotationally invariant, multi-scale
descriptors that capture the morphological content of a field at a
hierarchy of spatial scales and orientation interactions. They are
strictly more informative than two-point statistics for fields with
non-Gaussian structure and have been used as field-level summary
statistics in ISM, CMB, and galaxy-morphology inference contexts
\citep{CITATION_NEEDED_TBD}. \CGK{Suggested references to fill in:
Allys+19, Cheng+20, R\'egaldo-Saint Blancard+20/+23; please verify
ADS keys and substitute.}

In the long term, the goal is field-level inference: training a neural
posterior estimator directly on simulated maps, with wavelet scattering
coefficients (or learned summary statistics) serving as the
intermediate representation. This program collapses the distinction
between forward modeling and inference: each \TIGRESS\ snapshot
produces, via the pipeline of \autoref{sec:disc_synth}, the same
observables that \PHANGS\ measures, and the neural estimator learns the
posterior over $\thetaenv$ (and eventually $\thetaphys$) directly from
the resolved map. The companion paper takes the first step toward this
goal using aperture-level summary statistics; the present suite is sized
and structured so that the same simulations can be reused as the
training set as field-level summaries are added.

\subsection{Connection to the Companion Inference Paper}
\label{sec:disc_paperII}

The simulation suite presented here is designed from the outset to
serve as the numerical prior and simulator for the simulation-based
inference (\SBI) framework developed in the companion paper
\CGK{(\citealt{CITATION_NEEDED_TBD}; Paper~II, in preparation)}. The
sampling design of \autoref{sec:sampling}, mapped through the KDE of
the \PHANGS\ joint distribution, defines the implicit prior $p(\thetaenv)$
under which the inference is performed. Summary statistics extracted
from each \TIGRESS\ snapshot---$\Ssfr$, $\fmol$, $\PDE$, $\seff$, and
the decomposed feedback yields $\Yth$ and $\Ynonth$---constitute the
simulated observables that are matched against \PHANGS\ aperture
measurements. The progressive extensibility of the Sobol design ensures
that the training set can be enlarged without invalidating any inference
performed on earlier subsets, and the reservation of $\fmol$ and
$\Ssfr$ as validation quantities means that the inferred posteriors can
be cross-checked against quantities that did not enter the design. The
present paper completes the construction of the simulator; the
companion paper applies it.
